\chapter{Totale Beschränktheit und Kompaktheit}
\hypertarget{compactness.proofs}{}
\noindent\CompactnessMainLink{}\enspace\textperiodcentered\enspace\CompactnessReadingLink{}
\par\medskip

Die Definitionen und Hauptsätze stehen in Band~47. Hier werden ihre Beweise
ausgeführt. Nur die drei zusätzlich benötigten Hilfssätze erhalten eigene
Nummern mit dem Bandkennzeichen \(47\mathrm E\).
Es seien stets \((X,d)\) ein metrischer Raum und \(K\subseteq X\).
Wir schreiben \(B_r^K(x)=B_r^{d_K}(x)\) für die Kugel im Teilraum.
Für Punkte aus \(K\) gilt \(d_K(x,y)=d(x,y)\).

\section{Endliche Netze und getrennte Folgen}

\CompactnessProofTarget{MetricTotalBoundedSetBounded}
\noindent\textbf{Beweis zu \FormulaRefAuto{MetricTotalBoundedSetBounded}[thm].}
\begin{proof}
Für \(K=\varnothing\) gilt Beschränktheit nach Definition. Sei
\(K\ne\varnothing\). Wähle ein endliches \(1\)-Netz \(F\subseteq K\).
Da \(F\) die nichtleere Menge \(K\) überdeckt, ist \(F\ne\varnothing\).
Fixiere \(c_0\in F\). Es gibt eine reelle Zahl \(R>0\) mit
\(d(c,c_0)<R\) für alle \(c\in F\): Dies folgt durch endliche Induktion
über \(F\), indem man beim Hinzufügen eines Zentrums die bisherige
Schranke durch eine größere Schranke ersetzt. Der Induktionsschritt
verwendet nur die Vergleichbarkeit zweier reeller Zahlen.

Zu \(x\in K\) gibt es ein \(c\in F\) mit \(d(x,c)<1\). Daher
\[
 d(x,c_0)\leq d(x,c)+d(c,c_0)<1+R.
\]
Somit liegt \(K\) in der Kugel \(B_{1+R}^d(c_0)\), ist also beschränkt.
\end{proof}

\CompactnessProofTarget{MetricSeparatedSequenceFromNoFiniteNet}
\FormulaThmDeltaK[Eine fehlende endliche Überdeckung liefert eine getrennte Folge]{%
  \begin{aligned}[t]
  &\operatorname{Metrik}_{X}(d)\dsep K\subseteq X\dsep\varepsilon>0\dsep
    \neg\exists F\,\operatorname{Netz}_{d}(F,K,\varepsilon)
    \vdash\\[-2pt]
  &\quad\exists a\colon\mathbb N\to K\,
    \forall m,n\in\mathbb N\,
    \bigl(m\ne n\rightarrow\varepsilon\leq d(a_m,a_n)\bigr).
  \end{aligned}
}{MetricSeparatedSequenceFromNoFiniteNet}{%
  \DeltaRow{Mengen}{X\dsep K\dsep F\dsep\varepsilon\dsep m\dsep n}
  \DeltaRow{Funktionen}{d\dsep a}
}
\begin{proof}
\noindent\textbf{(i) Ein neuer Punkt außerhalb jeder endlichen Kugelfamilie.}
Für jede endliche Teilmenge \(F\subseteq K\) ist die Menge
\[
 R(F)=K\setminus\bigcup_{c\in F}B_\varepsilon^d(c)
\]
nichtleer. Andernfalls wäre \(F\) ein Netz im Sinne von
\FormulaRefAuto{MetricFiniteEpsilonNetDef}[def]. Insbesondere ist
\(K=R(\varnothing)\ne\varnothing\).

\medskip\noindent\textbf{(ii) Auswahl und Rekursion.}
Das Auswahlprinzip \FormulaRefAuto{Auswahlaxiom}[ax] liefert eine
Wahlfunktion \(s\) auf den nichtleeren Teilmengen von \(K\), mit
\(s(A)\in A\). Auf der Menge
\(\mathcal E=\{F\subseteq K\mid\Finite{F}\}\) ist deshalb
\[
 H(F)=F\cup\{s(R(F))\}
\]
eine wohldefinierte Abbildung \(\mathcal E\to\mathcal E\).
Dass der Wert endlich bleibt, folgt aus
\FormulaRefAuto{FiniteUnion}[thm] und der Endlichkeit einer Einermenge.
Der Rekursionssatz \FormulaRefAuto{RecursiveSequenceExistenceUnique}[thm]
liefert die Folge \(F_0=\varnothing\), \(F_{n+1}=H(F_n)\).
Setze \(a_n=s(R(F_n))\). Damit ist \(a\colon\mathbb N\to K\)
definiert; die Wahl der ganzen Folge ist nicht stillschweigend vorausgesetzt.

\medskip\noindent\textbf{(iii) Die Abstände bleiben getrennt.}
Induktion ergibt \(F_n=\{a_j\mid j<n\}\). Ist \(m<n\), so gehört
\(a_m\) zu \(F_n\), während \(a_n\) außerhalb jeder
\(\varepsilon\)-Kugel um einen Punkt von \(F_n\) liegt. Also
\(d(a_n,a_m)\geq\varepsilon\). Symmetrie liefert dieselbe Aussage für
\(n<m\), und damit die Behauptung.
\end{proof}

\Needspace{16\baselineskip}
\section{Das Cauchy-Teilfolgenkriterium}

\CompactnessProofTarget{MetricTotalBoundedCauchySubsequenceCriterion}
\noindent\textbf{Beweis zu
\FormulaRefAuto{MetricTotalBoundedCauchySubsequenceCriterion}[thm].}
\begin{proof}
\noindent\textbf{(i) Totale Beschränktheit liefert Cauchy-Teilfolgen.}
Sei \(K\) total beschränkt und \(a\colon\mathbb N\to K\) eine Folge.
Damit ist \(K\ne\varnothing\). Setze \(r_n=1/(n+1)\).
Nach \FormulaRefAuto{ReciprocalNaturalSequenceConvergesToZero}[thm]
gilt \(r_n\to0\). Für jedes \(n\) gibt es ein endliches
\(r_n\)-Netz \(F_n\). Wir wählen diese Netze mittels
\FormulaRefAuto{Auswahlaxiom}[ax] aus den nichtleeren Familien
\[
 \{F\subseteq K\mid\operatorname{Netz}_d(F,K,r_n)\}.
\]
Außerdem fixieren wir nach demselben Auswahlprinzip eine Wahlfunktion
\(s\) auf den nichtleeren Teilmengen von \(K\).

Für eine unendliche Indexmenge \(J\subseteq\mathbb N\) sei
\[
 C(n,J)=\{c\in F_n\mid
   \{j\in J\mid d(a_j,c)<r_n\}\text{ ist unendlich}\}.
\]
Diese Menge ist nichtleer. Denn die endlich vielen Mengen
\(\{j\in J\mid d(a_j,c)<r_n\}\) für \(c\in F_n\) überdecken \(J\).
Wären sie alle endlich, wäre \(J\) nach endlicher Induktion und
\FormulaRefAuto{FiniteUnion}[thm] endlich. Dieses Argument verlangt nicht,
dass die Kugeln disjunkt sind.

Mit \(J_0=\mathbb N\) definieren wir rekursiv
\[
 c_n=s(C(n,J_n)),\qquad
 J_{n+1}=\{j\in J_n\mid d(a_j,c_n)<r_n\}.
\]
Formal wendet man \FormulaRefAuto{RecursiveSequenceExistenceUnique}[thm]
auf den Zustand \((n,J)\) in
\(\mathbb N\times\{J\subseteq\mathbb N\mid\Infinite{J}\}\) und den
Schritt \((n,J)\mapsto(n+1,\{j\in J\mid d(a_j,s(C(n,J)))<r_n\})\) an.
Jede Menge \(J_n\) ist unendlich, und \(J_{n+1}\subseteq J_n\).

Jede unendliche Teilmenge von \(\mathbb N\) besitzt Elemente oberhalb
jeder vorgegebenen natürlichen Zahl: Eine beschränkte Teilmenge läge in
einem endlichen Anfangsabschnitt, wäre nach
\FormulaRefAuto{FiniteNatSegSubset}[thm] also endlich. Deshalb sind die
folgenden Minima durch
\FormulaRefAuto{NaturalWellOrderingPrinciple}[thm] wohldefiniert:
\[
 \varphi(0)=\min J_1,\qquad
 \varphi(n+1)=\min\{j\in J_{n+2}\mid\varphi(n)<j\}.
\]
Auch dies ist eine Rekursion auf Zuständen \((n,p)\); sie liefert eine
streng wachsende Teilfolgenabbildung mit \(\varphi(n)\in J_{n+1}\).

Sind \(m,\ell\geq n\), so gehören wegen der Verschachtelung beide
Indizes \(\varphi(m),\varphi(\ell)\) zu \(J_{n+1}\). Folglich
\[
 \begin{aligned}
 d(a_{\varphi(m)},a_{\varphi(\ell)})
 &\leq d(a_{\varphi(m)},c_n)+d(c_n,a_{\varphi(\ell)})\\
 &<2r_n.
 \end{aligned}
\]
Zu \(\eta>0\) wählen wir \(n\) so groß, dass \(r_n<\eta/2\).
Die letzte Abschätzung gilt für alle \(m,\ell\geq n\), also ist die
Teilfolge Cauchy. Da alle Werte in \(K\) liegen, gilt dies für \(d_K\)
ebenso wie für \(d\). Für \(K=\varnothing\) gibt es keine Folge
\(\mathbb N\to K\); die universelle Aussage ist dann ebenfalls wahr.

\medskip\noindent\textbf{(ii) Cauchy-Teilfolgen erzwingen totale Beschränktheit.}
Angenommen, jede Folge in \(K\) besitze eine Cauchy-Teilfolge, aber \(K\)
sei nicht total beschränkt. Nach
\FormulaRefAuto{MetricTotalBoundednessDef}[def] gibt es dann ein
\(\varepsilon>0\) ohne endliches \(\varepsilon\)-Netz.
\FormulaRefAuto{MetricSeparatedSequenceFromNoFiniteNet}[thm] liefert
eine Folge \(a\) mit \(d(a_m,a_n)\geq\varepsilon\) für \(m\ne n\).
Für jede streng wachsende \(\varphi\) und jeden Index \(N\) gilt daher
\[
 d(a_{\varphi(N)},a_{\varphi(N+1)})\geq\varepsilon.
\]
Keine Teilfolge kann die Cauchy-Bedingung zum Radius \(\varepsilon\)
erfüllen. Dies widerspricht der Voraussetzung.
\end{proof}

\section{Die Charakterisierung der Folgenkompaktheit}

\CompactnessProofTarget{MetricSequentiallyCompactSetTotallyBounded}
\FormulaThmDeltaK[Folgenkompakte Teilmengen sind total beschränkt]{%
  \operatorname{Metrik}_{X}(d)\dsep K\subseteq X\dsep
  \operatorname{Folgenkompakt}_{d}(K)
  \vdash\operatorname{Totalbeschr}_{d}(K)
}{MetricSequentiallyCompactSetTotallyBounded}{%
  \DeltaRow{Mengen}{X\dsep K}
  \DeltaRow{Funktionen}{d}
}
\begin{proof}
Jede Folge in \(K\) hat nach Voraussetzung eine Teilfolge, die gegen einen
Punkt von \(K\) konvergiert. Nach
\FormulaRefAuto{MetricSubspaceConvergenceAgree}[thm] konvergiert sie auch
bezüglich \(d_K\); nach
\FormulaRefAuto{MetricConvergentSequenceIsCauchy}[thm] ist sie dort Cauchy.
Das bereits bewiesene
\FormulaRefAuto{MetricTotalBoundedCauchySubsequenceCriterion}[thm]
liefert totale Beschränktheit. Die Aussage gilt auch für die leere Menge.
\end{proof}

\CompactnessProofTarget{MetricSequentialCompactnessCharacterization}
\noindent\textbf{Beweis zu
\FormulaRefAuto{MetricSequentialCompactnessCharacterization}[thm].}
\begin{proof}
\noindent\textbf{(i) Folgenkompaktheit liefert beide Bedingungen.}
Nach \FormulaRefAuto{MetricSequentiallyCompactSetComplete}[thm] ist der
Teilraum \((K,d_K)\) vollständig. Der gerade bewiesene Hilfssatz
\FormulaRefAuto{MetricSequentiallyCompactSetTotallyBounded}[thm]
liefert seine totale Beschränktheit.

\medskip\noindent\textbf{(ii) Vollständigkeit und totale Beschränktheit
liefern Folgenkompaktheit.}
Sei \(a\colon\mathbb N\to K\). Nach
\FormulaRefAuto{MetricTotalBoundedCauchySubsequenceCriterion}[thm]
existiert eine streng wachsende \(\varphi\), für die \(a\circ\varphi\)
bezüglich \(d_K\) Cauchy ist. Die Vollständigkeit von \((K,d_K)\)
liefert ein \(x\in K\) mit \(a_{\varphi(n)}\MetricTo{d_K}x\).
Nach \FormulaRefAuto{MetricSubspaceConvergenceAgree}[thm] ist dies
zugleich Konvergenz bezüglich \(d\). Damit ist
\FormulaRefAuto{MetricSequentialCompactnessDef}[def] erfüllt.
Für \(K=\varnothing\) sind die Folgenaussagen leer erfüllt und die
Netze leer; beide Richtungen bleiben gültig.
\end{proof}

\Needspace{16\baselineskip}
\section{Offene Überdeckungen}

\CompactnessProofTarget{MetricSequentiallyCompactLebesgueNumber}
\FormulaThmDeltaK[Ein gemeinsamer Kugelradius für eine offene Überdeckung]{%
  \begin{aligned}[t]
  &\operatorname{Metrik}_{X}(d)\dsep K\subseteq X\dsep
    \operatorname{Folgenkompakt}_{d}(K)\dsep
    \mathcal U\subseteq\mathcal P(K)\dsep\\[-2pt]
  &\quad\bigcup\mathcal U=K\dsep
    \forall U\in\mathcal U\;\operatorname{Offen}_{d_K}(U)
    \vdash\\[-2pt]
  &\quad\exists\lambda>0\;\forall x\in K\;\exists U\in\mathcal U\;
    B_\lambda^{d_K}(x)\subseteq U.
  \end{aligned}
}{MetricSequentiallyCompactLebesgueNumber}{%
  \DeltaRow{Mengen}{X\dsep K\dsep\mathcal U\dsep U\dsep x\dsep\lambda}
  \DeltaRow{Funktionen}{d\dsep d_K}
}
\begin{proof}
Für \(K=\varnothing\) erfüllt \(\lambda=1\) die Aussage.
Angenommen, für nichtleeres \(K\) existiere kein solcher Radius.
Für jedes \(n\in\mathbb N\) ist dann die Menge
\[
 A_n=\{x\in K\mid\forall U\in\mathcal U\;
                B_{1/(n+1)}^K(x)\not\subseteq U\}
\]
nichtleer. \FormulaRefAuto{Auswahlaxiom}[ax] liefert eine Folge
\(x_n\in A_n\). Die Folgenkompaktheit liefert eine streng wachsende
\(\varphi\) und ein \(x\in K\) mit \(x_{\varphi(j)}\MetricTo{d}x\).

Da \(\mathcal U\) überdeckt, gibt es ein \(U_0\in\mathcal U\) mit
\(x\in U_0\). Seine Offenheit liefert \(r>0\) mit
\(B_r^K(x)\subseteq U_0\). Für hinreichend großes \(j\) gilt zugleich
\[
 d(x_{\varphi(j)},x)<r/2,\qquad
 \frac1{\varphi(j)+1}<r/2.
\]
Die zweite Ungleichung folgt aus \(\varphi(j)\geq j\) und
\FormulaRefAuto{ReciprocalNaturalSequenceConvergesToZero}[thm].
Für \(y\in B_{1/(\varphi(j)+1)}^K(x_{\varphi(j)})\) ergibt sich
\[
 d(y,x)\leq d(y,x_{\varphi(j)})+d(x_{\varphi(j)},x)<r.
\]
Also ist diese ganze Kugel in \(B_r^K(x)\subseteq U_0\) enthalten.
Das widerspricht \(x_{\varphi(j)}\in A_{\varphi(j)}\).
\end{proof}

Der Hilfssatz ist die hier benötigte Kugelform des Lebesgue-Zahl-Lemmas.
Er benutzt nur Folgenkompaktheit und Offenheit, nicht schon die Existenz
einer endlichen Teilüberdeckung.

\CompactnessProofTarget{MetricCompactnessEquivalences}
\noindent\textbf{Beweis zu \FormulaRefAuto{MetricCompactnessEquivalences}[thm].}
\begin{proof}
\noindent\textbf{(i) Folgenkompaktheit impliziert Überdeckungskompaktheit.}
Sei \(\mathcal U\) eine offene Überdeckung von \(K\). Für
\(K=\varnothing\) genügt die leere Teilfamilie.
Für nichtleeres \(K\) liefert
\FormulaRefAuto{MetricSequentiallyCompactLebesgueNumber}[thm]
einen Radius \(\lambda>0\), so dass jede \(\lambda\)-Kugel im Teilraum
in einem Mitglied von \(\mathcal U\) liegt. Nach
\FormulaRefAuto{MetricSequentiallyCompactSetTotallyBounded}[thm]
besitzt \(K\) ein endliches \(\lambda\)-Netz \(F\).
Für jedes \(c\in F\) wählen wir ein \(U_c\in\mathcal U\) mit
\(B_\lambda^K(c)\subseteq U_c\). Diese endlich vielen Wahlen lassen
sich durch endliche Induktion treffen; auch das vorhandene Auswahlprinzip
wäre anwendbar. Das Bild \(\mathcal V=\{U_c\mid c\in F\}\) ist eine
endliche Teilfamilie von \(\mathcal U\), und
\[
 K=\bigcup_{c\in F}B_\lambda^K(c)
   \subseteq\bigcup\mathcal V\subseteq K.
\]
Somit ist \(\mathcal V\) eine endliche Teilüberdeckung.

\medskip\noindent\textbf{(ii) Überdeckungskompaktheit impliziert
Folgenkompaktheit.}
Sei \(K\) kompakt und \(a\colon\mathbb N\to K\) eine Folge.
Angenommen, sie habe keine in \(K\) konvergente Teilfolge.
Dann gibt es zu jedem \(x\in K\) ein \(r>0\), für das
\[
 I(x,r)=\{n\in\mathbb N\mid a_n\in B_r^K(x)\}
\]
endlich ist. Denn andernfalls gäbe es einen Punkt \(x\in K\), für den
alle diese Indexmengen unendlich wären. Man könnte dann rekursiv die
jeweils kleinsten Indizes wählen mit
\[
 \begin{gathered}
 \varphi(0)=\min I(x,1),\\
 \varphi(n+1)=\min\{j\in I(x,1/(n+2))\mid\varphi(n)<j\}.
 \end{gathered}
\]
Die Minima existieren wegen der Unbeschränktheit unendlicher Teilmengen
von \(\mathbb N\); Rekursion und Wohlordnung sind wie im
Cauchy-Teilfolgenbeweis anwendbar. Es gälte
\(d(a_{\varphi(n)},x)<1/(n+1)\), also
\(a_{\varphi(n)}\MetricTo{d}x\), ein Widerspruch.

Wir betrachten nun die Familie \emph{aller} Kugeln mit dieser Eigenschaft:
\[
 \mathcal W=\{B_r^K(x)\mid x\in K,\ r>0,\ \Finite{I(x,r)}\}.
\]
Sie ist eine offene Überdeckung von \(K\), weil jeder Punkt Mittelpunkt
mindestens einer solchen Kugel ist. Ihre Offenheit folgt aus
\FormulaRefAuto{MetricOpenBallsAreOpen}[thm] im Teilraum.
Eine gesonderte Auswahl von Radien ist für die Bildung dieser Familie
nicht nötig. Die Kompaktheit liefert eine endliche Teilüberdeckung
\(\mathcal V\subseteq\mathcal W\). Zu jedem \(V\in\mathcal V\)
ist \(\{n\in\mathbb N\mid a_n\in V\}\) endlich. Wegen der Überdeckung
ist jedoch
\[
 \mathbb N=\bigcup_{V\in\mathcal V}\{n\in\mathbb N\mid a_n\in V\}.
\]
Endliche Induktion und \FormulaRefAuto{FiniteUnion}[thm] würden
\(\mathbb N\) endlich machen, im Widerspruch zu
\FormulaRefAuto{NaturalNumbersInfinite}[thm]. Also besitzt jede Folge
eine in \(K\) konvergente Teilfolge. Für \(K=\varnothing\) ist die
Folgenaussage wiederum leer erfüllt.

\medskip\noindent\textbf{(iii) Die dritte Beschreibung.}
Die Schritte (i) und (ii) ergeben
\[
 \operatorname{Kompakt}_d(K)\leftrightarrow
 \operatorname{Folgenkompakt}_d(K).
\]
Zusammen mit \FormulaRefAuto{MetricSequentialCompactnessCharacterization}[thm]
folgt die Gleichwertigkeit mit Vollständigkeit und totaler Beschränktheit.
\end{proof}
